% sup2_03_S7b: Supplement part 2, step S7, second half: Cor. 2.3 re-proved with its exact scope; the classification.
% Input from sup2_00_intro.tex's section (through INPUTS_part2.tex).

\begin{lemma}[{\cite[Cor.~2.3]{Ba3}, with its exact scope}]\label{sup2:lem:S7cor}
Let $d=3$ and $G$ connected with at least one line.
\begin{enumerate}[label=\textup{(\alph*)},nosep,leftmargin=*]
\item If $G$ has a vertex of type (1.13)--(1.15), then $D(G)\ge1$.
\item If $G$ has an $R$-vertex \cite[(1.9)]{Ba3} or \cite[(1.11)]{Ba3} and $\bar n\ge3$, then $D(G)\ge\bar n-2\ge1$ (\cite{Ba3} takes $\bar n\ge12$, \cite[p.~421]{Ba3}).
\item If $G$ has a vertex \cite[(1.7)]{Ba3} with a finite counterterm, then: $E\ge2$ implies $D(G)\ge1$; $E=1$ implies $D(G)\ge0$, with $D(G)=0$ exactly for one graph $G_{\mathrm{tad}}$, the vertex $(1.7)_{\mathrm{fin}}$ and a vertex \cite[(1.8)]{Ba3} with $n=1$, $n'=0$ joined by two scalar lines, the $A'$ leg external; $E=0$ allows $D(G)=-1$ (the single $(1.7)_{\mathrm{fin}}$ vertex with its two legs contracted, $D=2-d$) and $D(G)=0$. So the third clause of \cite[Cor.~2.3]{Ba3} holds for $E\ge2$ and fails as printed for vacuum graphs and for $G_{\mathrm{tad}}$.
\item If $E\ge5$, then $D(G)\ge1$.
\end{enumerate}
\end{lemma}

\begin{proof}
Degrees lie in $\frac12\Z$ at $d=3$. Let $\cV_5$ be the set of (1.13)--(1.15) vertices and $\cV'$ the rest.

(a) By \cref{sup2:lem:S7vert}(b) at $d=3$, $D_G(v)=\frac52+\frac12e^\phi_v+\frac12n_{\mathrm{int}}(v)$ for $v\in\cV_5$, where $e^\phi_v=1$ if the $\phi'$ leg of $v$ is external and $0$ otherwise, and $D_G(v')\ge\frac12+\frac12e_{v'}$ for $v'\in\cV'$. Put $E'\defeq\sum_{\cV'}e_{v'}+\sum_{\cV_5}e^\phi_v$. Then $D(G)\ge-3+\frac52|\cV_5|+\frac12|\cV'|+\frac12E'+\frac12\sum n_{\mathrm{int}}$. If $|\cV_5|\ge2$, $D(G)\ge2$. Let $\cV_5=\{v\}$. Every vertex of $\cV'$ has $2$ or $4$ scalar legs and every scalar line uses two, while $v$ has one $\phi'$ leg; so the number of external scalar legs, which are counted in $E'$, is odd, and $E'\ge1$. If $|\cV'|\ge2$: $D(G)\ge-3+\frac52+1+\frac12=1$. If $\cV'=\emptyset$: all lines are self-loops at $v$, hence vector self-loops ($v$ has one scalar leg), so $n_{\mathrm{int}}\ge2$ and the $\phi'$ leg is external: $D(G)\ge-3+\frac52+\frac12+1=1$. If $\cV'=\{v'\}$: $D(G)=D_G(v)+D_G(v')-3$. If $D(v')\ge1$, then $D(G)\ge\frac52+1+\frac12E'-3\ge1$. If $D(v')=\frac12$, then $v'$ is \cite[(1.8)]{Ba3} with $n=1$, $n'=0$ (two scalar legs and one $A'$ leg; \cref{sup2:lem:S7vert}(a)). If the $\phi'$ leg of $v$ is external, $v$ and $v'$ are joined by a vector line only, so $n_{\mathrm{int}}\ge1$, both scalar legs of $v'$ are external, and $D(G)\ge\frac52+\frac12+\frac12+\frac12-3=1$. If the $\phi'$ leg of $v$ lies on the line to $v'$, the other scalar leg of $v'$ is external and its $A'$ leg is external ($e_{v'}=2$) or joins $v$ ($e_{v'}=1$, $n_{\mathrm{int}}=1$); in both cases $D(G)\ge\frac52+\frac12+\frac12+\frac12-3=1$.

(b) If $G$ has a vertex $w\in\cV_5$, then with $u$ the $R$-vertex $D(G)\ge D_G(w)+D_G(u)-3\ge\frac52+(\bar n+1)-3=\bar n+\frac12$, all other $D_G\ge\frac12$ being positive. Otherwise \cref{sup2:lem:S7id} with $\kappa_u\ge\bar n+1$ and all other terms $\ge0$: a tree has $E\ge2$ ($\sum(l_v-2)=E-2\ge0$), so $\ell+E\ge1$ and $D(G)\ge1-4+\bar n+1=\bar n-2$.

(c) If $G$ has a vertex $w\in\cV_5$, then with $u$ the $(1.7)_{\mathrm{fin}}$ vertex ($D_G(u)\ge D(u)=2$), $D(G)\ge\frac52+2-3=\frac32$, which implies every clause. Otherwise \cref{sup2:lem:S7id} gives $D(G)\ge\ell+E-2+\sum_{v\ne u}\kappa_v+\Delta_{\mathrm{ext}}$. $E\ge3$ gives $D(G)\ge1$. $E=2$: $\ell\ge1$ gives $D(G)\ge1$; if $\ell=0$, then $\sum(l_v-2)=0$, every vertex has two legs, every two-leg vertex has $\kappa\ge1$ ((1.7): $\ge1$; (1.8) with $n=0$: $\kappa=n'\ge1$; (1.10) with $n=0$: $\kappa\ge2$; $R$-vertices: $\bar n+1$), and $V\ge2$, so $D(G)\ge1$. $E=1$: a tree is impossible ($\sum(l_v-2)=-1$); $\ell\ge2$ gives $D(G)\ge1$; if $\ell=1$, then $\sum(l_v-2)=1$, one vertex has three legs and all others two, and $D(G)=\sum_{v\ne u}\kappa_v+\Delta_{\mathrm{ext}}=0$ forces $\kappa_v=0$ at every other vertex; no two-leg vertex has $\kappa=0$, and the only three-leg vertex with $\kappa=0$ is \cite[(1.8)]{Ba3} with $n=1$, $n'=0$. So $V=2$, $G=G_{\mathrm{tad}}$, its external leg being the undifferentiated $A'$ leg ($\Delta_{\mathrm{ext}}=0$; $E_s$ is even since every vertex has $2$ or $4$ scalar legs and every scalar line uses two, so the external leg is a vector leg). $E=0$: the one-vertex ring has $\ell=1$ and $D=1-4+2=-1$; the pair $(1.7)_{\mathrm{fin}}$, \cite[(1.6)]{Ba3} joined by three lines has $\ell=2$ and $D=0$.

(d) If $G$ has a vertex of type (1.13)--(1.15), (a) gives $D(G)\ge1$. Otherwise \cref{sup2:lem:S7id} with $\kappa\ge0$ (\cref{sup2:lem:S7vert}(c) for the $\delta m^2_k$ vertices) gives $D(G)\ge E-4\ge1$.
\end{proof}

\begin{lemma}[the two exceptions are immaterial]\label{sup2:lem:S7exc}
\textup{(a)} At background $\tilde B=0$ (after the gauge transformation removing $B_0$, \cite[p.~433]{Ba3}) every localised value of $G_{\mathrm{tad}}$, and of every graph containing $G_{\mathrm{tad}}$ as a subgraph, is $0$. \textup{(b)} Vacuum graphs never rely on \cite[Cor.~2.3]{Ba3}: by \cite[p.~432]{Ba3} ``all vacuum energy counterterms corresponding to these graphs are added'' to every vacuum class, whatever its degree.
\end{lemma}

\begin{proof}
(a) At $\tilde B=0$ the scalar propagator is a scalar kernel times $\mathbf 1_N$. The vertex \cite[(1.7)]{Ba3} is $\phi'\cdot\delta m^2\phi'$, where $\delta m^2$ is a matrix: \cite[(1.22)--(1.23), pp.~416--417]{Ba3} contain $q^2$ and $q^4$, so $\delta m^2$ is an even polynomial $P(q^2)$ in $q$, real symmetric and commuting with $q$; this holds for $\delta m^2_{\mathrm{fin}}$ and $\delta m^2_{K,k}$ alike. The vertex \cite[(1.8)]{Ba3} with $n+n'=1$ is $[D\phi'(b)\cdot q\phi'(b_-)]A'_b$. Contracting the scalar loop gives a scalar kernel times $\Tr(qP(q^2))=0$, the trace of an antisymmetric matrix times a symmetric one commuting with it. The same trace occurs in every graph containing $G_{\mathrm{tad}}$. (b) is the printed rule.
\end{proof}

At $\tilde B\ne0$, before the cube step, $G_{\mathrm{tad}}$ need not vanish; it is not used there, because the one-vector-leg classes are formed ``by graphs necessary to renormalize all divergent proper subgraphs'' \cite[p.~431]{Ba3} and the $\tilde B'$-expansion vertices carry $\kappa\ge1$, which makes $D>0$. By the same count, the statement of \cite[p.~431]{Ba3} that graphs with ``one external vector field leg \dots\ do not introduce any new divergences'' is also true only after the gauge transformation, by \cref{sup2:lem:S7exc}(a).

\begin{corollary}[the classification used in \cite{Ba3}, $d=3$]\label{sup2:cor:S7class}
Let $D(G)\le0$, $G$ not a vacuum graph and not $G_{\mathrm{tad}}$. Then $G$ has no vertex of type (1.13)--(1.15), no $R$-vertex and no finite-counterterm vertex, and $E\le4$; $E=4$ only for the graph \cite[(2.4)]{Ba3}; $E=3$ only with $(E_s,E_v)\in\{(0,3),(2,1)\}$, and then $\ell\le1$; $E=2$ only with $(E_s,E_v)\in\{(2,0),(0,2)\}$ and $V\le4$; $E=1$ only with $E_v=1$.
\end{corollary}

\begin{proof}
The first statement is \cref{sup2:lem:S7cor}. $E=4$: $D=\ell+\sum\kappa+\Delta_{\mathrm{ext}}\le0$ forces $\ell=0$, $\kappa=0$ everywhere and $\Delta_{\mathrm{ext}}=0$. A tree with four legs has $\sum(l_v-2)=2$, and a $\kappa=0$ vertex has $l_v\ge3$ ((1.8) with $n\ge1$, $n'=0$: $n+2$ legs; (1.10) with $n'=0$: $n\ge2$; (1.6): four legs); so $V\le2$, and $V=2$ because $G$ has a line. Both vertices are \cite[(1.8)]{Ba3} with $n=1$, $n'=0$. A vector line between them would leave both differentiated scalar legs external ($\Delta_{\mathrm{ext}}=2$), so they are joined by a scalar line carrying both differentiations: this is \cite[(2.4)]{Ba3}. $E=3$: $D=\ell-1+\dots\le0$ gives $\ell\le1$, and $E_s$ is even. $E=2$: \cref{sup2:lem:S7-217}(a) gives $D\ge\frac V2-2$, and $E_s$ is even. $E=1$: $E_s$ is even.
\end{proof}
